% !TEX program = xelatex
% Научная статья НИТУ «МИСиС»

\documentclass{misisarticle}

% polyglossia подменяет заголовок списка литературы; возвращаем «ЛИТЕРАТУРА» по шаблону
\gappto\captionsrussian{\renewcommand{\refname}{Литература}}

% --- Классификационные коды ---
\grnti{28.23.15, 28.23.37}
\asjc{1702, 1707}
\oecd{1.02, 1.03}

% --- Название ---
\title{Каскадный анализ калибровочных экспериментов сверхпроводящих кубитов: физическая аппроксимация с~эскалацией к~визуально-языковой модели}

% --- Авторы ---
\misisauthor{Р.~Я.~Наумов}{кафедра инженерной кибернетики НИТУ «МИСиС»}{Москва, Россия}{m2304542@edu.misis.ru}

% --- Аннотация и ключевые слова ---
\abstract{%
  Калибровка сверхпроводящих квантовых процессоров требует многократного анализа результатов типовых экспериментов: спектроскопии, осцилляций Раби, эксперимента Рамзи, измерения времени релаксации. Сегодня этот анализ выполняет инженер или большая визуально-языковая модель (VLM), читающая графики. Вызывать VLM для каждого эксперимента дорого и избыточно: большинство результатов рутинны и надёжно описываются физической моделью сигнала. В работе предлагается двухступенчатый каскад. На первой ступени сырые данные аппроксимируются моделью соответствующего эксперимента, а классификатор по признакам качества аппроксимации определяет исход (норма или тип аномалии) и рекомендуемое действие. Примеры, где уверенность классификатора ниже порога, передаются на вторую ступень~--- VLM, которая получает график, результаты аппроксимации и описание эксперимента. Для обучения и проверки разработан генератор синтетических экспериментов пяти типов с шестью классами исходов и известной разметкой; визуальная ступень дополнительно проверяется на открытом бенчмарке QCalEval. Качество оценивается по точности определения исхода и действия, ошибке извлечения параметров и доле примеров, переданных VLM, в сравнении с анализом только аппроксимацией и только VLM.
}

\keywords{%
  калибровка кубитов, сверхпроводящие кубиты, визуально-языковые модели, каскадный вывод, аппроксимация физической моделью, синтетические данные
}

\begin{document}

\maketitle

% ============================================================
\section{Введение}

Сверхпроводящие кубиты~--- одна из ведущих платформ для построения квантовых процессоров~\cite{krantz2019}. В России активно развивается направление флаксониевых кубитов: продемонстрированы высокоточные двухкубитные операции на флаксониумах с перестраиваемой связью~\cite{moskalenko2022}, а в 2026~г. в НИТУ «МИСиС» представлен прототип 16-кубитного квантового вычислителя на флаксониумах~\cite{misis2026}.

Чтобы кубиты выполняли операции с высокой точностью, их параметры необходимо регулярно калибровать: определять частоты кубитов и резонаторов считывания, амплитуды и длительности управляющих импульсов, времена релаксации и дефазировки~\cite{krantz2019}. Калибровка представляет собой последовательность экспериментов, результат каждого из которых нужно проанализировать и решить, принять ли полученный параметр или изменить условия измерения. Параметры дрейфуют во времени, поэтому процедура повторяется постоянно, а с ростом числа кубитов ручной анализ становится узким местом.

Методы искусственного интеллекта всё шире применяются на всех уровнях квантового вычислительного стека, и автоматическая калибровка рассматривается как ближайший путь к автономной работе устройств~\cite{alexeev2025}. Появились агентные системы на основе больших языковых моделей, которые самостоятельно ведут калибровку~\cite{cao2025kagents}, и визуально-языковые модели (VLM), интерпретирующие калибровочные графики~\cite{cao2026qcaleval, nvidia2026ising}. В 2026~г. с их помощью выполнена полностью автоматическая калибровка 20-кубитного процессора~\cite{qcr2026iqm}.

Однако вызов VLM с десятками миллиардов параметров для анализа каждого эксперимента избыточен: в реальной лаборатории доступны сырые данные, и большинство результатов рутинны~--- их надёжно описывает физическая модель сигнала, как это делается в классических пакетах анализа~\cite{kanazawa2023}. При этом классический анализ плохо распознаёт нестандартные исходы и не объясняет причину неудачи. В бенчмарке QCalEval~\cite{cao2026qcaleval} сравнивались только VLM; совместное использование физической аппроксимации и VLM не исследовалось.

Цель работы~--- снизить вычислительную стоимость автоматического анализа калибровочных экспериментов без потери точности за счёт каскадной схемы. Вклад работы:
\begin{itemize}
  \item каскад из двух ступеней: аппроксимация физической моделью с классификатором исхода по признакам качества аппроксимации и эскалация неуверенных случаев к VLM;
  \item генератор синтетических калибровочных экспериментов с известной разметкой исходов, параметров и рекомендуемых действий;
  \item протокол оценки, связывающий точность анализа с долей обращений к VLM.
\end{itemize}

Раздел~II содержит анализ существующих работ, раздел~III~--- постановку задачи, раздел~IV~--- описание метода и протокола проверки.

% ============================================================
\section{Анализ существующих работ}

\subsection{Классический анализ калибровочных экспериментов}

Стандартный подход к анализу~--- нелинейная аппроксимация измеренной кривой моделью эксперимента: лоренцевым контуром для спектроскопии, косинусом для осцилляций Раби, затухающей синусоидой для эксперимента Рамзи, экспонентой для измерения $T_1$~\cite{krantz2019}. Этот подход реализован, например, в пакете Qiskit Experiments~\cite{kanazawa2023}, где модуль анализа кривых выполняет аппроксимацию и оценивает её качество по приведённому $\chi^2$. Метод быстр и интерпретируем, но решение о качестве принимается по фиксированным порогам, а причина неудачной аппроксимации (сигнал вне диапазона, недостаточная дискретизация, выбросы) не определяется.

\subsection{Машинное обучение для калибровки}

Обзор~\cite{alexeev2025} показывает, что обучение с подкреплением, компьютерное зрение и байесовская оптимизация применяются для автоматической настройки, компенсации дрейфа и формирования импульсов. Так, обучение с подкреплением использовано для синтеза устойчивых к ошибкам гейтов на сверхпроводящем процессоре~\cite{baum2021}. Авторы обзора отмечают, что автоматическая калибровка пока недостаточно надёжна для работы без надзора человека, что делает важной интерпретируемость решений.

\subsection{Агентные системы и визуально-языковые модели}

В работе~\cite{cao2025kagents} предложены k-агенты: агенты на основе языковых моделей хранят знания лаборатории о доступных операциях и методах анализа, разбивают многошаговые процедуры на конечные автоматы и в замкнутом цикле калибруют сверхпроводящий процессор. Бенчмарк QCalEval~\cite{cao2026qcaleval} содержит 243 калибровочных графика из 22 семейств экспериментов с шестью типами вопросов: описание графика, классификация исхода, научное обоснование, оценка надёжности аппроксимации, извлечение параметров и диагностика. Лучшая модель общего назначения набрала в режиме zero-shot 72{,}3 балла, специализированная модель NVIDIA Ising Calibration~--- 74{,}7. Следующая версия этой модели содержит 31~млрд параметров~\cite{nvidia2026ising}. Эти работы используют графики как универсальное представление, но не задействуют сырые данные и не оценивают стоимость анализа.

Мультимодальные модели и методы компьютерного зрения применяются и в прикладных работах кафедры: для текстового поиска трёхмерных моделей~\cite{chetvertkov2026} и распознавания производственных дефектов~\cite{ereshchenko2023}. В качестве открытой VLM в данной работе рассматривается семейство Qwen2.5-VL~\cite{bai2025qwen}.

\subsection{Каскадный вывод и избирательная классификация}

Идея передавать сложные запросы дорогой модели, а простые обрабатывать дешёвой, развита в работе FrugalGPT~\cite{chen2023frugal}, где каскад языковых моделей снижает стоимость при сохранении качества. Теоретическую основу правила эскалации даёт избирательная классификация~\cite{geifman2017}: модель отказывается от ответа, если уверенность ниже порога, а качество описывается кривой «риск~--- покрытие». Применительно к анализу калибровочных экспериментов такие схемы не исследовались.

% ============================================================
\section{Постановка задачи}

Калибровочный эксперимент описывается кортежем $e = (k, \mathbf{x}, \mathbf{y}, c)$, где $k \in K$~--- тип эксперимента, $\mathbf{x} \in \mathbb{R}^n$~--- сетка сканирования, $\mathbf{y} \in \mathbb{R}^n$~--- результаты измерений, $c$~--- контекст (параметры сканирования, описание эксперимента). Рассматриваются типы
\[
  K = \{\text{рез.}, \ \text{спектр.}, \ \text{Раби}, \ \text{Рамзи}, \ T_1\}.
\]

Требуется построить отображение $f(e) = (o, \hat{\theta}, a)$, где $o \in O$~--- исход эксперимента из шести классов (табл.~\ref{tab:classes}), $\hat{\theta}$~--- оценка физических параметров (частота, амплитуда $\pi$-импульса, $T_2^*$, $T_1$), $a = g(o)$~--- рекомендуемое действие.

Имеются две модели: дешёвая $F_1$ со стоимостью $C_1$ и дорогая $F_2$ со стоимостью $C_2 \gg C_1$. Модель $F_1$ возвращает распределение $p_1(o \mid e)$. Каскад с порогом $\tau$ определяется как
\begin{equation}
  f_\tau(e) =
  \begin{cases}
    F_1(e), & \max_{o} p_1(o \mid e) \geq \tau, \\
    F_2(e), & \text{иначе}.
  \end{cases}
  \label{eq:cascade}
\end{equation}
Ожидаемая стоимость анализа одного эксперимента
\begin{equation}
  C(\tau) = C_1 + \rho(\tau)\, C_2,
  \label{eq:cost}
\end{equation}
где $\rho(\tau)$~--- доля эскалаций. Задача состоит в выборе порога, минимизирующего долю эскалаций при ограничении на точность:
\begin{equation}
  \tau^* = \arg\min_{\tau} \rho(\tau) \quad \text{при} \quad \mathrm{Acc}(f_\tau) \geq \mathrm{Acc}(F_2) - \varepsilon,
  \label{eq:objective}
\end{equation}
где $\varepsilon$~--- допустимая потеря точности.

\begin{table}[htbp]
  \caption{Классы исходов и рекомендуемые действия}
  \label{tab:classes}
  \centering
  \small
  \begin{tabular}{p{0.27\columnwidth}p{0.6\columnwidth}}
    \toprule
    \textbf{Исход} & \textbf{Действие} \\
    \toprule
    Норма & принять параметр \\
    Сильный шум & увеличить число усреднений \\
    Нет сигнала & проверить настройку, расширить поиск \\
    Вне диапазона & сдвинуть или расширить окно сканирования \\
    Редкая сетка & уменьшить шаг сканирования \\
    Скачок, выбросы & повторить измерение \\
    \bottomrule
  \end{tabular}
\end{table}

% ============================================================
\section{Описание метода}

\subsection{Генератор синтетических экспериментов}

Для обучения первой ступени и контролируемой проверки каскада нужны эксперименты с известным исходом и истинными параметрами. Открытый бенчмарк QCalEval содержит только изображения графиков без сырых данных~\cite{cao2026qcaleval}, поэтому разработан генератор синтетических экспериментов. Сигналы задаются моделями:
\begin{align}
  y_{\text{L}}(x) &= b + A\frac{(\gamma/2)^2}{(x - x_0)^2 + (\gamma/2)^2}, \label{eq:lorentz}\\
  y_{\text{R}}(x) &= b + \frac{A}{2}\left(1 - \cos\frac{2\pi x}{P}\right), \label{eq:rabi}\\
  y_{\text{Rm}}(x) &= b + \frac{A}{2}\left(1 + e^{-x/T_2^*}\cos 2\pi\Delta x\right), \label{eq:ramsey}\\
  y_{T_1}(x) &= b + A e^{-x/T_1}, \label{eq:t1}
\end{align}
где \eqref{eq:lorentz} описывает спектроскопию резонатора (провал, $A<0$) и кубита (пик, $A>0$), \eqref{eq:rabi}~--- осцилляции Раби с амплитудой $\pi$-импульса $P/2$, \eqref{eq:ramsey}~--- эксперимент Рамзи с расстройкой $\Delta$, \eqref{eq:t1}~--- релаксацию. К сигналу добавляется гауссов шум.

Аномальные исходы моделируются изменением условий измерения: сильный шум~--- увеличением дисперсии шума; отсутствие сигнала~--- нулевой амплитудой; выход за диапазон~--- смещением резонанса за пределы окна или слишком коротким либо длинным окном по времени; редкая сетка~--- числом точек меньше двух на период; скачок~--- сдвигом уровня или выбросами, имитирующими флуктуации дефектов. Доли классов задаются так, чтобы преобладали рутинные исходы (55\,\% нормы). Каждый пример содержит сырые данные, график и разметку: тип, исход, истинные параметры и действие. Примеры приведены на рис.~\ref{fig:examples}.

\begin{figure*}[htbp]
  \centering
  \includegraphics[width=0.92\textwidth]{figures/fig_examples.png}
  \caption{Примеры синтетических калибровочных экспериментов: нормальные исходы и аномалии.}
  \label{fig:examples}
\end{figure*}

\subsection{Ступень 1: аппроксимация и классификатор исхода}

Для эксперимента типа $k$ параметры модели находятся нелинейным методом наименьших квадратов (алгоритм Левенберга~--- Марквардта). Начальное приближение вычисляется по данным: положение экстремума для лоренцева контура, доминирующая частота спектра Фурье для осцилляций, уровень плато для экспоненты.

По результату аппроксимации вычисляется вектор признаков:
\begin{itemize}
  \item коэффициент детерминации $R^2$ и приведённый $\chi^2$, нормированный на оценку шума по разностям соседних точек;
  \item относительные стандартные отклонения параметров из ковариационной матрицы;
  \item положение найденного параметра относительно окна сканирования и число периодов в окне;
  \item число точек на период или на ширину линии;
  \item отношение амплитуды сигнала к шуму;
  \item статистики остатков: максимальный стандартизованный остаток и автокорреляция первого порядка, чувствительная к скачкам.
\end{itemize}

На признаках и типе эксперимента обучается градиентный бустинг, предсказывающий исход $o$. Вероятности калибруются на валидационной выборке, после чего максимальная вероятность используется как уверенность в правиле~\eqref{eq:cascade}. Для уверенных примеров ступень возвращает исход, параметры аппроксимации и действие $g(o)$.

\subsection{Ступень 2: визуально-языковая модель}

Неуверенные примеры передаются VLM. В отличие от постановки QCalEval~\cite{cao2026qcaleval}, модель получает не только график, но и результаты первой ступени: найденные параметры, значения признаков качества и распределение $p_1(o \mid e)$, а также текстовое описание эксперимента. Ответ запрашивается в структурированном виде~--- исход, параметр, действие и краткое объяснение для инженера~--- и проверяется на соответствие схеме. В качестве VLM рассматриваются открытая модель Qwen2.5-VL~\cite{bai2025qwen} и специализированная NVIDIA Ising Calibration~\cite{nvidia2026ising}. Схема каскада показана на рис.~\ref{fig:cascade}.

\begin{figure}[htbp]
  \centering
  \includegraphics[width=0.85\columnwidth]{figures/fig_cascade.png}
  \caption{Схема каскадного анализа калибровочного эксперимента.}
  \label{fig:cascade}
\end{figure}

\subsection{Выбор порога эскалации}

Порог $\tau$ выбирается на валидационной выборке по задаче~\eqref{eq:objective}. Для каждого значения $\tau$ строятся доля эскалаций $\rho(\tau)$ и точность каскада; полученная кривая аналогична кривой «риск~--- покрытие» избирательной классификации~\cite{geifman2017}. Стоимости $C_1$ и $C_2$ в~\eqref{eq:cost} измеряются как время анализа одного эксперимента на одинаковом оборудовании.

\subsection{Протокол проверки}

Проверка проводится на двух наборах данных:
\begin{itemize}
  \item синтетический набор из генератора (не менее 3000 экспериментов) с разбиением на обучающую, валидационную и тестовую части по независимым начальным значениям генератора;
  \item бенчмарк QCalEval~\cite{cao2026qcaleval} для оценки визуальной ступени на графиках с реального оборудования.
\end{itemize}

Сравниваются три конфигурации: только ступень~1, только VLM и каскад. Метрики: точность и macro-F1 классификации исхода, точность выбора действия, относительная ошибка извлечения параметров для нормальных исходов, доля эскалаций $\rho$ и среднее время анализа. Отдельно оценивается устойчивость к уровню шума и к типу эксперимента, исключённому из обучения первой ступени.

\subsection{Ограничения и дальнейшая работа}

Синтетические данные не воспроизводят всех особенностей реальных устройств, в частности специфики флаксониевых кубитов. Следующий этап~--- проверка каскада на архивных калибровочных данных реального процессора и встраивание каскада в замкнутый цикл автоматической калибровки.

% ============================================================
\begin{thebibliography}{99}

\bibitem{krantz2019} Krantz P., Kjaergaard M., Yan F., Orlando T.~P., Gustavsson S., Oliver W.~D. A quantum engineer's guide to superconducting qubits. \textit{Applied Physics Reviews}, vol.~6, no.~2, art.~021318, 2019. DOI: 10.1063/1.5089550.

\bibitem{moskalenko2022} Moskalenko I.~N. et al. High fidelity two-qubit gates on fluxoniums using a tunable coupler. \textit{npj Quantum Information}, vol.~8, art.~130, 2022. DOI: 10.1038/s41534-022-00644-x.

\bibitem{misis2026} 16-кубитный квантовый компьютер на флаксониумах создали в МИСИС [Электронный ресурс]. 2026. URL: \url{https://www1.ru/news/2026/02/22/16-kubitnyi-kvantovyi-kompiuter-na-flaksoniumax-sozdali-v-misis.amp.html}.

\bibitem{alexeev2025} Alexeev Y., Farag M.~H., Patti T.~L., Wolf M.~E., Ares N. et al. Artificial intelligence for quantum computing. \textit{Nature Communications}, 2025. DOI: 10.1038/s41467-025-65836-3.

\bibitem{cao2025kagents} Cao S., Zhang Z., Alghadeer M., Fasciati S.~D., Piscitelli M., Bakr M. et al. Agents for self-driving laboratories applied to quantum computing. \textit{Patterns}, 2025. DOI: 10.1016/j.patter.2025.101372.

\bibitem{cao2026qcaleval} Cao S., Zhang Z. et al. QCalEval: Benchmarking vision-language models for quantum calibration plot understanding. arXiv:2604.25884, 2026.

\bibitem{nvidia2026ising} NVIDIA Ising enables fully automated quantum computer calibration with enhanced in-context learning [Электронный ресурс]. NVIDIA Technical Blog, 2026. URL: \url{https://developer.nvidia.com/blog/nvidia-ising-enables-fully-automated-quantum-computer-calibration-with-enhanced-in-context-learning}.

\bibitem{qcr2026iqm} IQM and NVIDIA automate 20-qubit chip calibration via NVIDIA Ising vision language models and QCalEval benchmark [Электронный ресурс]. Quantum Computing Report, 2026. URL: \url{https://quantumcomputingreport.com/iqm-and-nvidia-automate-20-qubit-chip-calibration-via-nvidia-ising-vision-language-models-and-qcaleval-benchmark/}.

\bibitem{kanazawa2023} Kanazawa N., Egger D.~J., Ben-Haim Y., Zhang H., Shanks W.~E., Aleksandrowicz G., Wood C.~J. Qiskit Experiments: A Python package to characterize and calibrate quantum computers. \textit{Journal of Open Source Software}, vol.~8, no.~84, art.~5329, 2023. DOI: 10.21105/joss.05329.

\bibitem{baum2021} Baum Y. et al. Experimental deep reinforcement learning for error-robust gate-set design on a superconducting quantum computer. \textit{PRX Quantum}, vol.~2, art.~040324, 2021.

\bibitem{chetvertkov2026} Четвертков Е.~С. Сравнительный анализ мультимодальных архитектур в задаче текстового поиска 3D-моделей LEGO // Искусственный интеллект в промышленных, коммерческих, медицинских и финансовых приложениях: сб. ст. науч.-техн. семинара. Вып.~6. М.: НИТУ «МИСИС», 2026. С.~141--146.

\bibitem{ereshchenko2023} Ерещенко А.~Г. Исследование распознавания производственных дефектов на стальных поверхностях при помощи компьютерного зрения // Искусственный интеллект в промышленных, коммерческих, медицинских и финансовых приложениях: сб. ст. науч.-техн. семинара студентов кафедры «Инженерной кибернетики». Вып.~1. М.: НИТУ «МИСИС», 2023. С.~139--144.

\bibitem{bai2025qwen} Bai S. et al. Qwen2.5-VL technical report. arXiv:2502.13923, 2025.

\bibitem{chen2023frugal} Chen L., Zaharia M., Zou J. FrugalGPT: How to use large language models while reducing cost and improving performance. arXiv:2305.05176, 2023.

\bibitem{geifman2017} Geifman Y., El-Yaniv R. Selective classification for deep neural networks. \textit{Advances in Neural Information Processing Systems}, vol.~30, 2017.

\end{thebibliography}

\end{document}
